% ============================================================
%  front matter -- series page, title page, copyright, dedication, preface
%  (printed pages i-xii of the original)
% ============================================================

% ---- half title (printed p. i) ----
\thispagestyle{empty}
\vspace*{\fill}
\begin{center}
\large\itshape Introduction to Fourier Analysis

on Euclidean Spaces
\end{center}
\vfill
\cleardoublepage

% ---- series page (printed p. ii) ----
\thispagestyle{empty}
\begin{center}
\bfseries PRINCETON MATHEMATICAL SERIES
\end{center}
\vspace{2ex}
\begin{center}
Editors: WU-CHUNG HSIANG, JOHN MILNOR, and ELIAS M. STEIN
\end{center}
\vspace{2ex}

{\parindent=0pt
\noindent\hangindent=1.8em\hangafter=1
1. The Classical Groups. By HERMANN WEYL

\noindent\hangindent=1.8em\hangafter=1
3. An Introduction to Differential Geometry, By LUTHER PFAHLER EISENHART

\noindent\hangindent=1.8em\hangafter=1
4. Dimension Theory, By W. HUREWICZ and W. WALLMAN

\noindent\hangindent=1.8em\hangafter=1
6. The Laplace Transform, By D. V. WIDDER

\noindent\hangindent=1.8em\hangafter=1
7. Integration. By EDWARD J. McSHANE

\noindent\hangindent=1.8em\hangafter=1
8. Theory of Lie Groups: 1, By C. CHEVALLEY

\noindent\hangindent=1.8em\hangafter=1
9. Mathematical Methods of Statistics, By HARALD CRAM\'ER

\noindent\hangindent=1.8em\hangafter=1
10. Several Complex Variables, By S. BOCHNER and W. T. MARTIN

\noindent\hangindent=1.8em\hangafter=1
11. Introduction to Topology, By S. LEFSHETZ

\noindent\hangindent=1.8em\hangafter=1
12. Algebraic Geometry and Topology, edited By R. H. Fox, D. C. SPENCER, and A. W. TUCKER

\noindent\hangindent=1.8em\hangafter=1
14. The Topology of Fibre Bundles, By NORMAN STEENROD

\noindent\hangindent=1.8em\hangafter=1
15. Foundations of Algebraic Topology, By SAMUEL EILENBERG and NORMAN STEENROD

\noindent\hangindent=1.8em\hangafter=1
16. Functionals of Finite Riemann Surfaces, By MENAHEM SCHIFFER and DONALD C. SPENCER.

\noindent\hangindent=1.8em\hangafter=1
17. Introduction to Mathematical Logic, Vol. 1, By Alonzo Church

\noindent\hangindent=1.8em\hangafter=1
19. Homological Algebra, By H. Cartan and S. Eilenberg

\noindent\hangindent=1.8em\hangafter=1
20. The Convolution Transform, By I. I. Hirschman and D. V. Widder

\noindent\hangindent=1.8em\hangafter=1
21. Geometric Integration Theory, By H. WHITNEY

\noindent\hangindent=1.8em\hangafter=1
22. Qualitative Theory of Differential Equations, By V. V. NEMYTSKII and V. V. STEPANOV

\noindent\hangindent=1.8em\hangafter=1
23. Topological Analysis, By GORDON T. WHYBURN (revised 1964)

\noindent\hangindent=1.8em\hangafter=1
24. Analytic Functions, By AHLFORS, BEHNKE and GRAUERT, BERS, et al.

\noindent\hangindent=1.8em\hangafter=1
25. Continuous Geometry, By JOHN VON NEUMANN

\noindent\hangindent=1.8em\hangafter=1
26. Riemann Surfaces, by L. AHLFORS and L. SARIO

\noindent\hangindent=1.8em\hangafter=1
27. Differential and Combinatorial Topology, edited By S. S. CAIRNS

\noindent\hangindent=1.8em\hangafter=1
28. Convex Analysis, By R. T. ROCKAFELLAR

\noindent\hangindent=1.8em\hangafter=1
29. Global Analysis, edited By D. C. SPENCER and S. IYANAGA

\noindent\hangindent=1.8em\hangafter=1
30. Singular Integrals and Differentiability Properties of Functions, By E. M. STEIN

\noindent\hangindent=1.8em\hangafter=1
31. Problems in Analysis, edited By R. C. GUNNING

\noindent\hangindent=1.8em\hangafter=1
32. Introduction to Fourier Analysis on Euclidean Spaces, By E. M. STEIN and G. WEISS
\par}

\cleardoublepage

% ---- title page (printed p. iii) ----
\thispagestyle{empty}
\vspace*{2cm}
\begin{center}
\Large Introduction to {\Large FOURIER ANALYSIS}

\Large ON EUCLIDEAN SPACES
\end{center}
\vspace{4ex}
\begin{center}
\large By Elias M. Stein \& Guido Weiss
\end{center}
\vspace*{\fill}
\begin{center}
PRINCETON, NEW JERSEY

PRINCETON UNIVERSITY PRESS
\end{center}
\vspace{2ex}
\cleardoublepage

% ---- copyright page (printed p. iv) ----
\thispagestyle{empty}
\vspace*{3cm}
\begin{center}
Copyright \textcopyright\ 1971 by Princeton University Press

L.C. Card: 73-106394

ISBN: 0-691-08078-X

AMS 1970: 31B05, 31B25, 32A07, 33A45, 42A02, 42A18, 42A40, 42A68, 44A25, 46E30, 46E35

All Rights Reserved

Second Printing, with corrections, 1975

Sixth Printing, 1990
\end{center}
\vspace{2ex}
\begin{center}
13 15 14 12
\end{center}
\vspace{2ex}
\begin{center}
\small Princeton University Press books are printed on acid-free paper and
meet the guidelines for permanence and durability of the Committee on
Production Guidelines for Book Longevity of the Council on Library Resources
\end{center}
\vspace{2ex}
\begin{center}
\itshape This book has been composed in Times New Roman type

Printed in the United States of America

\url{http://pup.princeton.edu}

ISBN-13: 978-0-691-08078-9 (cloth)

ISBN-10: 0-691-08078-X (cloth)
\end{center}
\cleardoublepage

% ---- dedication (printed p. v) ----
\thispagestyle{empty}
\vspace*{\fill}
\begin{center}
THIS BOOK IS DEDICATED TO

\vspace{2ex}
\large Antoni Zygmund

\vspace{3ex}
IN APPRECIATION FOR HIS FRIENDSHIP,

HIS TEACHING, AND THE INSPIRATION

HE GAVE US
\end{center}
\vspace*{\fill}
\cleardoublepage

% ---- preface (printed pp. vii-viii) ----
\chapter*{Preface}
\addcontentsline{toc}{chapter}{Preface}
\markboth{PREFACE}{}

This book is designed to be an introduction to harmonic analysis in Euclidean
spaces. The subject has seen a considerable flowering during the past twenty
years. We have not tried to cover all phases of this development. Rather, our
chief concern was to illustrate various methods used in this aspect of Fourier
analysis that exploit the structure of Euclidean spaces. In particular, we try
to show the role played by the action of translations, dilations, and rotations.
Another concern, not independent of this chief one, is to motivate the study of
harmonic analysis on more general spaces having an analogous structure (such as
arises in symmetric spaces). It is our feeling that the study of Fourier
analysis in that context and, also, in other general settings, is more
meaningful once the special Euclidean case is understood.

Because of these concerns we have not included several topics that are usually
presented in more general treatments of harmonic analysis. For example, results
centering around the Wiener Tauberian theorem, which holds in the case of
general locally compact Abelian groups, do not involve the special features of
Euclidean spaces. Stated very briefly, our selection of topics was motivated by
showing how real variable and complex variable methods extend from the
one-dimensional to the many-dimensional case.

We require that the readers have mastered the material that is usually covered
in a course in the theory of integration and in the theory of functions of a
complex variable. We also believe that this book is much more meaningful to
someone who has some previous acquaintance with harmonic analysis. If the reader
has no such previous experience, he might find it profitable first to look at
the expository article ``Harmonic Analysis'' (see Guido Weiss \cite{Weiss3}).

The idea of writing this book first occurred to us when we presented some of
these topics in graduate courses given during the academic year 1958--59. Since
that time each of us has lectured on this subject at various times and places.
We are thankful to our colleagues and students who helped us clarify our ideas
and organize this material.

More specifically, one of us (Weiss) taught a course on this subject at
Washington University jointly with Mitchell Taibleson during the academic year
1963--64. We are grateful for his contribution to the organization of this course
and his continued interest in our effort. The other author gave courses at the
University of Chicago (1961--62) and Princeton (1963--65) that dealt with the
topics of this book. Our task of writing the book was considerably facilitated
by having the lecture notes, which were prepared by R. Askey, N. J. Weiss, and
D. Levine. It is a pleasure to express our appreciation to them. Ronald R.
Coifman was a constant help throughout the writing of this book. His many
suggestions are incorporated in much of the presentation. Miguel de Guzm\'an,
J. R. Hattemer, Stephen Wainger, and N. J. Weiss read the entire manuscript. We
are grateful for their corrections and comments. We are also indebted to
I. I. Hirschman Jr., who read part of the manuscript and made several useful
suggestions, and to Mrs. A. Bonami, J. L. Clerc, L. J. Dickson, S. S. Gelbart,
and S. Zucker who helped us correct the proof sheets.

\vspace{2ex}
\noindent
\hangindent=0pt
\begin{tabular}{@{}l}
Princeton, N.J.\\
St. Louis, Mo.
\end{tabular}

\cleardoublepage

% ---- half title before Chapter I (printed p. xi) ----
\thispagestyle{empty}
\vspace*{\fill}
\begin{center}
\large\itshape Introduction to Fourier Analysis

on Euclidean Spaces
\end{center}
\vfill
\cleardoublepage
